\documentclass[review,11pt,3p,authoryear,nopreprintline]{elsarticle}

\usepackage[T1]{fontenc}
\usepackage{amsmath,amssymb,amsthm}
\usepackage{mathrsfs}
\usepackage{bm}
\usepackage{xcolor}
\usepackage[colorlinks=true,citecolor=blue,linkcolor=blue,urlcolor=blue]{hyperref}

\newtheorem{theorem}{Theorem}
\newtheorem{proposition}{Proposition}
\newtheorem{corollary}{Corollary}
\newtheorem{lemma}{Lemma}
\newdefinition{example}{Example}
\newdefinition{definition}{Definition}
\newdefinition{remark}{Remark}
\newdefinition{property}{Property}
\newdefinition{condition}{Condition}

\begin{document}

\begin{frontmatter}
  \title{Failure Probability Bounds with Uncertain Surrogate Residual Distributions}

  \author[inst1]{Renlong Wang}
  \ead{rw106@rice.edu}
  \affiliation[inst1]{organization={Department of Civil and Environmental Engineering, Rice University}, city={Houston}, postcode={TX 77005}, country={USA}}

  \begin{abstract}
  \end{abstract}

  \begin{keyword}
  \end{keyword}
\end{frontmatter}

\section{Introduction}

\section{Problem Description}
\label{sec-problem}

Let $\bm X\sim\mathbb P$ be a random input taking values in $\mathcal X\subseteq\mathbb R^n$, and let $g:\mathcal X\to\mathbb R$ be a measurable limit-state function. The failure probability is defined as
\[
  p_f:=\Pr\{g(\bm X)\le0\}=\mathbb E_{\bm X\sim\mathbb P}\left[\mathbf 1_{\{g(\bm X)\le0\}}\right]=\int_{\mathcal X}\mathbf 1_{\{g(\bm x)\le0\}}\,\mathbb P(\mathrm d\bm x).
\]
For fixed $g$, this functional is linear in $\mathbb P$.

Evaluating $g$ at an input $\bm x$ can be computationally expensive, particularly when each evaluation requires a high-fidelity simulation. A common approach is to replace $g$ with a computationally efficient surrogate model. Let
\[
  \mathcal D^{\mathrm{tr}}:=\left\{\left(\bm x_i^{\mathrm{tr}},g(\bm x_i^{\mathrm{tr}})\right)\right\}_{i=1}^{N^{\mathrm{tr}}}
\]
denote the training dataset used to construct a measurable surrogate $\hat g:=\hat g_{\mathcal D^{\mathrm{tr}}}$. All model fitting and selection are completed using the training data before the test data introduced below are used. Throughout the analysis, we condition on $\mathcal D^{\mathrm{tr}}$, treat $\hat g$ as fixed, and suppress this dependence in the notation.

Define the signed residual function
\[
  \xi(\bm x):=g(\bm x)-\hat g(\bm x),\qquad \forall \bm x\in\mathcal X.
\]
For a new input $\bm X\sim\mathbb P$ independent of the training data, the residual $\xi^\star:=\xi(\bm X)$ has distribution $\mathbb Q^\star$ under the fixed surrogate $\hat g$. Let $\Xi\subseteq\mathbb R$ be a known Borel set satisfying $\mathbb Q^\star(\Xi)=1$, and let $\mathcal Q(\Xi)$ denote the set of probability measures on $\mathbb R$ supported on $\Xi$. We take $\Xi=\mathbb R$ when no smaller support can be justified.

For a candidate residual distribution $\mathbb Q\in\mathcal Q(\Xi)$, let $\Pi$ denote a joint distribution of $(\bm X,\xi)$ on $\mathcal X\times\Xi$ with input marginal $\mathbb P = \Pi_{\bm X}$ and residual marginal $\mathbb Q = \Pi_{\xi}$. The corresponding surrogate-based failure probability is
\begin{equation}
  p_{\hat f}(\Pi):=\Pr_\Pi\{\hat g(\bm X)+\xi\le0\}=\int_{\mathcal X\times\Xi}\mathbf 1_{\{\hat g(\bm x)+\xi\le0\}}\,\Pi(\mathrm d\bm x,\mathrm d\xi).
  \label{eq-surrogate-failure-probability}
\end{equation}
For fixed $\hat g$, this functional is linear in $\Pi$. In particular, the true joint distribution $\Pi^\star$ satisfies $p_{\hat f}(\Pi^\star)=p_f$.

The surrogate-based failure probability depends on both the residual distribution and its dependence on the input. This dependence is particularly relevant for inputs whose surrogate predictions lie close to the failure threshold, where even small residual errors can change the failure classification. In complex engineering systems, however, the available information may be insufficient to specify the residual distribution reliably. We therefore construct an ambiguity set $\mathscr Q$ of plausible residual distributions intended to contain $\mathbb Q^\star$ with high confidence. For a given residual ambiguity set $\mathscr Q$, define the joint ambiguity set $\mathscr P(\mathscr Q)$ by
\[
  \mathscr P(\mathscr Q) :=\left\{\Pi\in\mathcal P(\mathcal X\times\Xi):\Pi_{\bm X}=\mathbb P,\ \Pi_\xi\in\mathscr Q\right\},
\]
where $\mathcal P(\mathcal X\times\Xi)$ is the set of probability measures on $\mathcal X\times\Xi$. This set allows all dependence structures consistent with the specified marginals. The upper and lower failure probability bounds are
\begin{equation}
  \overline p_{\hat f}:=\sup_{\Pi\in\mathscr P(\mathscr Q)}p_{\hat f}(\Pi),
  \label{eq-upper-bound}
\end{equation}
and
\begin{equation}
  \underline p_{\hat f}:=\inf_{\Pi\in\mathscr P(\mathscr Q)}p_{\hat f}(\Pi),
  \label{eq-lower-bound}
\end{equation}
respectively. If $\mathbb Q^\star\in\mathscr Q$, then $\Pi^\star\in\mathscr P(\mathscr Q)$ and hence $\underline p_{\hat f}\le p_f\le\overline p_{\hat f}$.

The main tasks are therefore:
\begin{enumerate}
  \item \textbf{Ambiguity set design and statistical calibration.} Construct a residual ambiguity set and establish its coverage guarantee under explicit sampling conditions.
  \item \textbf{Tractable reformulations.} Derive exact and computationally tractable reformulations of the infinite-dimensional optimization problems defining the failure probability bounds.
\end{enumerate}


\section{Ambiguity Set Design}
\label{sec-ambiguity-set}

In the surrogate formulation, failure occurs when $\xi\le-\hat g(\bm X)$. The residual cumulative distribution function (CDF) therefore describes the marginal probability mass relevant to each failure threshold. Its tail accuracy matters when estimating small failure probabilities. We use the D\"umbgen--Wellner (DW) construction \citep{L.J23} to control deviations from the empirical CDF simultaneously over all thresholds, with tolerances that vary across probability levels.

We first introduce a local discrepancy that accounts for the probability level. Let $q$ and $p$ denote an empirical probability and a candidate probability, respectively. Following \citet{B.J79} and the CDF confidence band construction of \citet{O95}, we compare them using the Bernoulli Kullback--Leibler divergence
\begin{equation}
  \phi(q,p):=q\log\frac{q}{p}+(1-q)\log\frac{1-q}{1-p},\qquad \forall q\in[0,1],\, p\in(0,1),
  \label{eq-bernoulli-kl}
\end{equation}
where logarithms are natural and a term with zero numerator is interpreted as zero. For $q$ close to an interior $p$, the approximation $\phi(q,p)\approx(q-p)^2/[2p(1-p)]$ scales the deviation by the Bernoulli variance. Thus, the same absolute deviation receives a larger weight near the probability boundaries. The full divergence also covers zero and full empirical counts through $\phi(0,p)=-\log(1-p)$ and $\phi(1,p)=-\log p$. At the boundaries of the candidate probability, we set $\phi(0,0)=\phi(1,1)=0$, $\phi(q,0)=+\infty$ for $q>0$, and $\phi(q,1)=+\infty$ for $q<1$.

To compare deviations simultaneously across probability levels, we use the additive correction of \citet{L.J23}. For $v\in(0,1)$, define
\begin{equation}
  \rho(v):=\log\log\frac{e}{4v(1-v)},\qquad \omega(v):=\log\!\left(1+\rho(v)^2\right).
  \label{eq-dw-correction-functions}
\end{equation}
We fix the divergence parameter $s=1$, corresponding to \eqref{eq-bernoulli-kl}, and the correction parameter $\nu=1$ before sampling. The resulting correction is $\psi(v):=\rho(v)+\omega(v)$, with $\psi(0)=\psi(1)=+\infty$. For $q,p\in[0,1]$, define its bivariate extension
\begin{equation}
  \psi(q,p):=\min_{v\in[\min\{q,p\},\,\max\{q,p\}]}\psi(v).
  \label{eq-dw-bivariate-correction}
\end{equation}
The functions $\rho$ and $\omega$ vanish at $v=1/2$ and increase toward the probability boundaries. Taking the minimum over the interval between $q$ and $p$ makes the correction large only when both probabilities lie in the same tail. This correction balances fluctuations across probability scales in the simultaneous CDF constraint.

We now construct the empirical CDF from a test dataset
\[
  \mathcal D^{\mathrm{te}}:=\left\{\left(\bm x_j^{\mathrm{te}},g(\bm x_j^{\mathrm{te}})\right)\right\}_{j=1}^{N^{\mathrm{te}}}.
\]
The test inputs $\bm x_1^{\mathrm{te}},\ldots,\bm x_{N^{\mathrm{te}}}^{\mathrm{te}}$ are drawn independently from $\mathbb P$ and are independent of $\mathcal D^{\mathrm{tr}}$. We fix $N^{\mathrm{te}}\ge1$ before sampling and use the test data to calibrate the residual ambiguity set. The fitted surrogate remains fixed during and after calibration.

Define the signed test residuals
\[
  \xi_j^{\mathrm{te}}:=\xi(\bm x_j^{\mathrm{te}})=g(\bm x_j^{\mathrm{te}})-\hat g(\bm x_j^{\mathrm{te}}),\qquad j=1,\ldots,N^{\mathrm{te}}.
\]
Conditional on $\mathcal D^{\mathrm{tr}}$, these residuals are independent and identically distributed according to $\mathbb Q^\star$. Their empirical CDF is
\begin{equation}
  \widehat F^{\mathrm{te}}(t):=\frac{1}{N^{\mathrm{te}}}\sum_{j=1}^{N^{\mathrm{te}}}\mathbf 1_{\{\xi_j^{\mathrm{te}}\le t\}},\qquad \forall t\in\mathbb R.
  \label{eq-empirical-cdf}
\end{equation}
For $\mathbb Q\in\mathcal Q(\Xi)$, write $F_{\mathbb Q}(t):=\mathbb Q(({-\infty},t])$. At threshold $t$, $q=\widehat F^{\mathrm{te}}(t)$ is the average of the Bernoulli indicators in \eqref{eq-empirical-cdf}, and $p=F_{\mathbb Q}(t)$ is the candidate probability. The quantity $N^{\mathrm{te}}\phi(q,p)$ is then the Bernoulli log-likelihood difference between the empirical and candidate probabilities. Combining this local discrepancy with the multiscale correction gives $N^{\mathrm{te}}\phi(q,p)-\psi(q,p)$.

At the probability boundaries, the corrected discrepancy is $-\infty$ for $q=p\in\{0,1\}$ and $+\infty$ for $p=0<q$ or $p=1>q$. These conventions extend the constraint to every residual threshold.

\begin{definition}[DW-based CDF ambiguity set]
  \label{definition-dw-ambiguity-set}
  Given the empirical CDF $\widehat F^{\mathrm{te}}$ and a threshold $\theta\ge0$, define
  \begin{equation} 
    \label{eq-cdf-ambiguity-set}
      \mathscr Q^{\mathrm{DW}}(\theta):=\left\{\mathbb Q\in\mathcal Q(\Xi):\ 
      \sup_{t\in\mathbb R}\left[ N^{\mathrm{te}}\phi\!\left(\widehat F^{\mathrm{te}}(t),F_{\mathbb Q}(t)\right) 
       {}-\psi\!\left(\widehat F^{\mathrm{te}}(t),F_{\mathbb Q}(t)\right)\right]\le\theta\right\}. 
  \end{equation} 
\end{definition}

We calibrate $\theta$ against a distribution-free reference law. Following \citet{L.J23}, let $V_1,\ldots,V_{N^{\mathrm{te}}}$ be independent standard uniform variables, independent of the training and test data. Their empirical CDF is
\begin{equation}
  G_{N^{\mathrm{te}}}(u):=\frac{1}{N^{\mathrm{te}}}\sum_{j=1}^{N^{\mathrm{te}}}\mathbf 1_{\{V_j\le u\}},\qquad \forall u\in[0,1].
  \label{eq-dw-reference-cdf}
\end{equation}
Define the uniform reference statistic
\begin{equation}
  T_{N^{\mathrm{te}}}^{\mathrm U}:=\sup_{u\in(0,1)}\left[N^{\mathrm{te}}\phi\!\left(G_{N^{\mathrm{te}}}(u),u\right)-\psi\!\left(G_{N^{\mathrm{te}}}(u),u\right)\right].
  \label{eq-dw-reference-statistic}
\end{equation}
Here $\Pr_{\mathrm U}$ denotes probability under the uniform reference law. The next proposition uses a quantile of this reference statistic to calibrate the residual ambiguity set.

\begin{proposition}[Finite-sample $\theta$ calibration]
  \label{proposition-01}
  Under the sampling conditions stated above, let $\delta\in(0,1)$ and define
  \begin{equation}
    \theta_{N^{\mathrm{te}}}(\delta):=\inf\left\{z\ge0:\Pr_{\mathrm U}\left\{T_{N^{\mathrm{te}}}^{\mathrm U}\le z\right\}\ge1-\delta\right\}.
    \label{eq-dw-calibrated-threshold}
  \end{equation}

  Then, for every $\theta\ge\theta_{N^{\mathrm{te}}}(\delta)$, the DW-based CDF ambiguity set satisfies
  \begin{equation}
    \Pr_{\mathcal D^{\mathrm{te}}}\left\{\mathbb Q^\star\in\mathscr Q^{\mathrm{DW}}(\theta)\,\middle|\,\mathcal D^{\mathrm{tr}}\right\}\ge1-\delta.
    \label{eq-cdf-coverage}
  \end{equation}
\end{proposition}
\begin{proof}
  The reference statistic is finite almost surely, and its distribution function is right-continuous. Hence \eqref{eq-dw-calibrated-threshold} gives
  \[
    \Pr_{\mathrm U}\left\{T_{N^{\mathrm{te}}}^{\mathrm U}>\theta_{N^{\mathrm{te}}}(\delta)\right\}\le\delta.
  \]
  Conditioning on $\mathcal D^{\mathrm{tr}}$ fixes $\hat g$ and $\mathbb Q^\star$. Define the generalized quantile function $F_{\mathbb Q^\star}^{-1}(u):=\inf\{t\in\mathbb R:F_{\mathbb Q^\star}(t)\ge u\}$ for all $u\in(0,1)$. The sample $F_{\mathbb Q^\star}^{-1}(V_j)$ for $j=1,\ldots,N^{\mathrm{te}}$, has the same conditional joint law as the test residuals. Thus, given $\mathcal D^{\mathrm{tr}}$,
  \[
    \left(\widehat F^{\mathrm{te}}(t)\right)_{t\in\mathbb R}\overset{\mathrm d}{=}\left(G_{N^{\mathrm{te}}}\!\left(F_{\mathbb Q^\star}(t)\right)\right)_{t\in\mathbb R}.
  \]
  For $0<F_{\mathbb Q^\star}(t)<1$, the discrepancy at $\mathbb Q^\star$ is evaluated at a probability level included in \eqref{eq-dw-reference-statistic}. When $F_{\mathbb Q^\star}(t)$ is zero or one, the empirical CDF equals that boundary almost surely, giving discrepancy $-\infty$. The supremum in \eqref{eq-cdf-ambiguity-set}, evaluated at $\mathbb Q^\star$, is therefore bounded by $T_{N^{\mathrm{te}}}^{\mathrm U}$ under this coupling. Since $\mathbb Q^\star(\Xi)=1$, we obtain
  \[
    \Pr_{\mathcal D^{\mathrm{te}}}\left\{\mathbb Q^\star\notin\mathscr Q^{\mathrm{DW}}\!\left(\theta_{N^{\mathrm{te}}}(\delta)\right)\,\middle|\,\mathcal D^{\mathrm{tr}}\right\}
    \le\Pr_{\mathrm U}\left\{T_{N^{\mathrm{te}}}^{\mathrm U}>\theta_{N^{\mathrm{te}}}(\delta)\right\}\le\delta.
  \]
  For every test sample, $\theta\ge\theta_{N^{\mathrm{te}}}(\delta)$ implies $\mathscr Q^{\mathrm{DW}}\!\left(\theta_{N^{\mathrm{te}}}(\delta)\right)\subseteq\mathscr Q^{\mathrm{DW}}(\theta)$. This inclusion proves \eqref{eq-cdf-coverage} for every such $\theta$.
\end{proof}

For every fixed $N^{\mathrm{te}}\ge1$, the proposition applies to arbitrary residual distributions, including laws with atoms or infinite moments. With $\phi$ and $\psi$ fixed as above, the reference distribution depends only on $N^{\mathrm{te}}$, so $\theta_{N^{\mathrm{te}}}(\delta)$ can be computed offline for each prescribed confidence level $1-\delta$. \citet[Supplement, Sections~S.6--S.7]{L.J23} describe finite-sample numerical calibration using probabilities that uniform order statistics lie within specified bounds. Choosing $\theta$ as a conservative upper bound on the quantile in \eqref{eq-dw-calibrated-threshold} preserves \eqref{eq-cdf-coverage}. If calibration is performed by simulation, quantile estimation error must be accounted for in the confidence guarantee. A threshold chosen from the test data also preserves \eqref{eq-cdf-coverage}, provided $\theta\ge\theta_{N^{\mathrm{te}}}(\delta)$ almost surely, by the same nesting argument.

The following theorem transfers residual-distribution coverage to simultaneous bounds on the failure probability under the joint model. It applies to any residual ambiguity set satisfying the stated coverage condition.

\begin{theorem}[Finite-sample failure probability bounds]
  \label{theorem-01}
  Let $\delta\in(0,1)$, and suppose a test-data-dependent residual ambiguity set $\mathscr Q$ satisfies
  \[
    \Pr_{\mathcal D^{\mathrm{te}}}\left\{\mathbb Q^\star\in\mathscr Q\,\middle|\,\mathcal D^{\mathrm{tr}}\right\}\ge1-\delta.
  \]
  Then the bounds defined in \eqref{eq-upper-bound} and \eqref{eq-lower-bound} satisfy
  \begin{equation}
    \Pr_{\mathcal D^{\mathrm{te}}}\left\{\underline p_{\hat f}\le p_f\le\overline p_{\hat f}\,\middle|\,\mathcal D^{\mathrm{tr}}\right\}\ge1-\delta.
    \label{eq-failure-probability-coverage}
  \end{equation}
  In particular, this guarantee holds for $\mathscr Q=\mathscr Q^{\mathrm{DW}}(\theta)$ whenever $\theta\ge\theta_{N^{\mathrm{te}}}(\delta)$, with the calibrated threshold defined in \eqref{eq-dw-calibrated-threshold}.
\end{theorem}
\begin{proof}
  On $\mathbb Q^\star\in\mathscr Q$, the true joint distribution $\Pi^\star$ belongs to $\mathscr P(\mathscr Q)$. Since $p_{\hat f}(\Pi^\star)=p_f$, the lower and upper bounds satisfy
  \[
    \underline p_{\hat f}\le p_{\hat f}(\Pi^\star)=p_f\le\overline p_{\hat f}.
  \]
  Taking probabilities yields
  \[
    \Pr_{\mathcal D^{\mathrm{te}}}\left\{\underline p_{\hat f}\le p_f\le\overline p_{\hat f}\,\middle|\,\mathcal D^{\mathrm{tr}}\right\}\ge\Pr_{\mathcal D^{\mathrm{te}}}\left\{\mathbb Q^\star\in\mathscr Q\,\middle|\,\mathcal D^{\mathrm{tr}}\right\}\ge1-\delta,
  \]
  which proves Theorem~\ref{theorem-01}.
\end{proof}

The event $\mathbb Q^\star\in\mathscr Q$ gives both inequalities, so $1-\delta$ is the confidence level for the entire interval. In the DW set, $\theta$ controls the corrected likelihood discrepancy. The corresponding CDF tolerances vary by probability level and can differ above and below the empirical CDF. For a fixed sample size, higher confidence requires a nondecreasing calibrated threshold $\theta_{N^{\mathrm{te}}}(\delta)$.

For a fixed empirical CDF, $0\le\theta_1\le\theta_2$ implies $\mathscr Q^{\mathrm{DW}}(\theta_1)\subseteq\mathscr Q^{\mathrm{DW}}(\theta_2)$. The lower bound is therefore nonincreasing in $\theta$, and the upper bound is nondecreasing. The interval reflects finite-sample uncertainty in the residual law and uncertainty in its dependence on the input. Because the joint ambiguity set allows all couplings with the specified marginals, a gap can persist even when the residual marginal is known exactly.

\section{Tractable Reformulation}
\label{sec-reformulation}

% Fix the training data, the test residuals, and a threshold $\theta\ge0$. We first express $\mathscr Q^{\mathrm{DW}}(\theta)$ through finitely many linear probability constraints and derive the Lagrangian duals of the failure probability problems. For the default support $\Xi=\mathbb R$, an exact aggregation then yields finite linear programs. Their dual structure gives closed-form bounds under the original input distribution $\mathbb P$. We assume that all test residuals belong to $\Xi$, as holds almost surely under the sampling conditions in Section~\ref{sec-ambiguity-set}.

% \subsection{Rank endpoint constraints}
% \label{subsec-rank-constraints}

% For $j=0,\ldots,N^{\mathrm{te}}$, let $q_j:=j/N^{\mathrm{te}}$ and define
% \begin{equation}
%   \mathcal I_j(\theta):=
%   \left\{p\in[0,1]:
%   N^{\mathrm{te}}\phi(q_j,p)-\psi(q_j,p)\le\theta\right\}.
%   \label{eq-dw-rank-acceptance}
% \end{equation}
% The corrected discrepancy decreases as $p$ approaches $q_j$ from below and increases as $p$ moves above $q_j$. Its acceptance set is therefore an interval,
% \begin{equation}
%   \mathcal I_j(\theta)=[\alpha_j(\theta),\beta_j(\theta)],
%   \qquad
%   \alpha_j(\theta):=\inf\mathcal I_j(\theta),\quad
%   \beta_j(\theta):=\sup\mathcal I_j(\theta).
%   \label{eq-dw-rank-endpoints}
% \end{equation}
% These are the rank endpoints obtained by inversion of the DW statistic \citep{L.J23}. They are nondecreasing in $j$ and satisfy
% \[
%   0\le\alpha_j(\theta)\le q_j\le\beta_j(\theta)\le1,
%   \qquad
%   \alpha_0(\theta)=0,\quad
%   \beta_{N^{\mathrm{te}}}(\theta)=1.
% \]
% For the symmetric choice $s=\nu=1$, $\alpha_j(\theta)=1-\beta_{N^{\mathrm{te}}-j}(\theta)$. The endpoints can be computed by one-dimensional root finding.

% Let $r_1<\cdots<r_K$ be the distinct test residuals, and set
% \begin{equation}
%   c_0:=0,\qquad
%   c_k:=\sum_{j=1}^{N^{\mathrm{te}}}\mathbf1_{\{\xi_j^{\mathrm{te}}\le r_k\}},
%   \qquad
%   u_k:=\beta_{c_{k-1}}(\theta),\quad
%   \ell_k:=\alpha_{c_k}(\theta).
%   \label{eq-dw-counts-and-levels}
% \end{equation}
% Thus $c_K=N^{\mathrm{te}}$; repeated residuals retain their multiplicities in $c_k$.

% \begin{proposition}[Finite rank constraints]
%   \label{prop-finite-rank-constraints}
%   The DW ambiguity set has the equivalent representation
%   \begin{equation}
%     \mathscr Q^{\mathrm{DW}}(\theta)=
%     \left\{\mathbb Q\in\mathcal Q(\Xi):
%     \begin{array}{ll}
%       \mathbb Q((-\infty,r_k))\le u_k,\\
%       \mathbb Q((-\infty,r_k])\ge\ell_k,
%     \end{array}
%     \quad k=1,\ldots,K\right\}.
%     \label{eq-finite-rank-constraints}
%   \end{equation}
% \end{proposition}
% \begin{proof}
%   For $k=1,\ldots,K-1$, the empirical CDF on $[r_k,r_{k+1})$ equals $c_k/N^{\mathrm{te}}$. Inverting the DW constraint gives
%   $\alpha_{c_k}(\theta)\le F_{\mathbb Q}(t)\le\beta_{c_k}(\theta)$ throughout this interval. By monotonicity of $F_{\mathbb Q}$, its lower constraint is equivalent to
%   $F_{\mathbb Q}(r_k)\ge\alpha_{c_k}(\theta)$, while its upper constraint is equivalent to
%   $F_{\mathbb Q}(r_{k+1}^-)\le\beta_{c_k}(\theta)$. Applying the same argument to the two tail intervals, using $\alpha_0(\theta)=0$ and $\beta_{N^{\mathrm{te}}}(\theta)=1$, gives exactly \eqref{eq-finite-rank-constraints}.
% \end{proof}

% The open upper half-line and closed lower half-line in \eqref{eq-finite-rank-constraints} account for possible atoms at the observed residuals.

% \subsection{Measure programs and Lagrangian duals}
% \label{subsec-measure-duals}

% Let $\mathcal M_+(\mathcal X\times\Xi)$ denote the finite nonnegative measures on $\mathcal X\times\Xi$. Substituting \eqref{eq-finite-rank-constraints} into \eqref{eq-upper-bound} gives the measure linear program
% \begin{equation}
%   \begin{aligned}
%     \overline p_{\hat f}
%     =\sup_{\Pi\in\mathcal M_+(\mathcal X\times\Xi)}\quad&
%     \int_{\mathcal X\times\Xi}
%     \mathbf1_{\{\hat g(\bm x)+\xi\le0\}}\,\Pi(\mathrm d\bm x,\mathrm d\xi)\\
%     \text{s.t.}\quad&
%     \Pi(A\times\Xi)=\mathbb P(A),
%     &&\forall A\in\mathcal B(\mathcal X),\\
%     &\int_{\mathcal X\times\Xi}
%     \mathbf1_{\{\xi<r_k\}}\,\Pi(\mathrm d\bm x,\mathrm d\xi)\le u_k,
%     &&k=1,\ldots,K,\\
%     &\int_{\mathcal X\times\Xi}
%     \mathbf1_{\{\xi\le r_k\}}\,\Pi(\mathrm d\bm x,\mathrm d\xi)\ge\ell_k,
%     &&k=1,\ldots,K.
%   \end{aligned}
%   \label{eq-measure-upper}
% \end{equation}
% The corresponding lower problem is
% \begin{equation}
%   \begin{aligned}
%     \underline p_{\hat f}
%     =\inf_{\Pi\in\mathcal M_+(\mathcal X\times\Xi)}\quad&
%     \int_{\mathcal X\times\Xi}
%     \mathbf1_{\{\hat g(\bm x)+\xi\le0\}}\,\Pi(\mathrm d\bm x,\mathrm d\xi)\\
%     \text{s.t.}\quad&
%     \Pi(A\times\Xi)=\mathbb P(A),
%     &&\forall A\in\mathcal B(\mathcal X),\\
%     &\int_{\mathcal X\times\Xi}
%     \mathbf1_{\{\xi<r_k\}}\,\Pi(\mathrm d\bm x,\mathrm d\xi)\le u_k,
%     &&k=1,\ldots,K,\\
%     &\int_{\mathcal X\times\Xi}
%     \mathbf1_{\{\xi\le r_k\}}\,\Pi(\mathrm d\bm x,\mathrm d\xi)\ge\ell_k,
%     &&k=1,\ldots,K.
%   \end{aligned}
%   \label{eq-measure-lower}
% \end{equation}
% The marginal equality implies $\Pi(\mathcal X\times\Xi)=1$. Both programs are feasible: the product of $\mathbb P$ and the empirical residual distribution satisfies all constraints.

% Introduce multipliers $\lambda,\mu\in\mathbb R_+^K$ and a bounded measurable function $\varphi:\mathcal X\to\mathbb R$. For the upper problem, the Lagrangian is
% \begin{equation}
%   \begin{aligned}
%     \mathcal L^+={}&
%     \int_{\mathcal X\times\Xi}\mathbf1_{\{\hat g(\bm x)+\xi\le0\}}\,\Pi(\mathrm d\bm x,\mathrm d\xi)
%     +\int_{\mathcal X}\varphi(\bm x)\,\mathbb P(\mathrm d\bm x)
%     -\int_{\mathcal X\times\Xi}\varphi(\bm x)\,\Pi(\mathrm d\bm x,\mathrm d\xi)\\
%     &+\sum_{k=1}^K\lambda_k
%     \left[u_k-\int_{\mathcal X\times\Xi}\mathbf1_{\{\xi<r_k\}}\,\Pi(\mathrm d\bm x,\mathrm d\xi)\right]\\
%     &+\sum_{k=1}^K\mu_k
%     \left[\int_{\mathcal X\times\Xi}\mathbf1_{\{\xi\le r_k\}}\,\Pi(\mathrm d\bm x,\mathrm d\xi)-\ell_k\right].
%   \end{aligned}
%   \label{eq-upper-lagrangian}
% \end{equation}
% Define the step function
% \begin{equation}
%   \eta_{\lambda,\mu}(\xi):=
%   \sum_{k=1}^K\lambda_k\mathbf1_{\{\xi<r_k\}}
%   -\sum_{k=1}^K\mu_k\mathbf1_{\{\xi\le r_k\}}.
%   \label{eq-dual-step-function}
% \end{equation}
% Collecting terms in \eqref{eq-upper-lagrangian} yields
% \[
%   \mathcal L^+
%   =
%   \mathbb E_{\mathbb P}[\varphi(\bm X)]
%   +\lambda^\top u-\mu^\top\ell
%   +\int_{\mathcal X\times\Xi}
%   \left[\mathbf1_{\{\hat g(\bm x)+\xi\le0\}}-\varphi(\bm x)
%   -\eta_{\lambda,\mu}(\xi)\right]\,\Pi(\mathrm d\bm x,\mathrm d\xi).
% \]
% Its supremum over $\Pi\ge0$ is finite exactly when the integrand is nonpositive at every point. A positive value at any point would allow an arbitrarily large point mass there. The upper dual is therefore
% \begin{equation}
%   \begin{aligned}
%     \overline d_\theta:=\inf_{\lambda,\mu\ge0,\varphi}\quad&
%     \mathbb E_{\mathbb P}[\varphi(\bm X)]
%     +\lambda^\top u-\mu^\top\ell\\
%     \text{s.t.}\quad&
%     \varphi(\bm x)+\eta_{\lambda,\mu}(\xi)
%     \ge\mathbf1_{\{\hat g(\bm x)+\xi\le0\}},
%     \quad\forall(\bm x,\xi)\in\mathcal X\times\Xi.
%   \end{aligned}
%   \label{eq-upper-measure-dual}
% \end{equation}

% For the lower problem, use $\lambda_k$ on the constraint
% $\int\mathbf1_{\{\xi<r_k\}}\,\mathrm d\Pi-u_k\le0$ and $\mu_k$ on
% $\ell_k-\int\mathbf1_{\{\xi\le r_k\}}\,\mathrm d\Pi\le0$. Its Lagrangian becomes
% \[
%   \mathcal L^-
%   =
%   \mathbb E_{\mathbb P}[\varphi(\bm X)]
%   -\lambda^\top u+\mu^\top\ell
%   +\int_{\mathcal X\times\Xi}
%   \left[\mathbf1_{\{\hat g(\bm x)+\xi\le0\}}-\varphi(\bm x)
%   +\eta_{\lambda,\mu}(\xi)\right]\,\Pi(\mathrm d\bm x,\mathrm d\xi).
% \]
% Taking its infimum over $\Pi\ge0$ requires a nonnegative integrand and gives
% \begin{equation}
%   \begin{aligned}
%     \underline d_\theta:=\sup_{\lambda,\mu\ge0,\varphi}\quad&
%     \mathbb E_{\mathbb P}[\varphi(\bm X)]
%     -\lambda^\top u+\mu^\top\ell\\
%     \text{s.t.}\quad&
%     \varphi(\bm x)-\eta_{\lambda,\mu}(\xi)
%     \le\mathbf1_{\{\hat g(\bm x)+\xi\le0\}},
%     \quad\forall(\bm x,\xi)\in\mathcal X\times\Xi.
%   \end{aligned}
%   \label{eq-lower-measure-dual}
% \end{equation}
% Weak duality gives
% $\underline d_\theta\le\underline p_{\hat f}\le\overline p_{\hat f}\le\overline d_\theta$.
% For fixed multipliers, the optimal upper potential is the pointwise supremum of
% $\mathbf1_{\{\hat g(\bm x)+\xi\le0\}}-\eta_{\lambda,\mu}(\xi)$ over $\xi\in\Xi$.
% The optimal lower potential is the corresponding infimum of
% $\mathbf1_{\{\hat g(\bm x)+\xi\le0\}}+\eta_{\lambda,\mu}(\xi)$.

% \subsection{Finite linear programs and strong duality}
% \label{subsec-finite-duality}

% We now take $\Xi=\mathbb R$. Set $J:=2K+1$, $r_0:=-\infty$, and $r_{K+1}:=+\infty$. Partition the residual axis into
% \begin{equation}
%   C_{2k}:=\{r_k\}\quad(k=1,\ldots,K),
%   \qquad
%   C_{2k+1}:=(r_k,r_{k+1})\quad(k=0,\ldots,K).
%   \label{eq-residual-cells}
% \end{equation}
% Define the input threshold and its cells by
% \begin{equation}
%   T(\bm x):=-\hat g(\bm x),\qquad
%   A_i:=\{\bm x:T(\bm x)\in C_i\},\qquad
%   w_i:=\mathbb P(A_i),\quad i=1,\ldots,J.
%   \label{eq-input-cells}
% \end{equation}
% On $A_i\times C_j$, the largest and smallest possible failure indicators are
% \begin{equation}
%   \overline h_{ij}:=\mathbf1_{\{j\le i\}},
%   \qquad
%   \underline h_{ij}:=\mathbf1_{\{j<i\}}
%   +\mathbf1_{\{j=i,\ i\ {\rm even}\}}.
%   \label{eq-cell-costs}
% \end{equation}
% The diagonal singleton cells always fail because equality is included in the failure event. On a diagonal open interval, a residual can be selected on either side of $T(\bm x)$.

% For a joint measure, let $p_{ij}:=\Pi(A_i\times C_j)$. The feasible cell masses form the polytope
% \begin{equation}
%   \mathcal A_\theta:=
%   \left\{p\in\mathbb R_+^{J\times J}:
%   \begin{array}{ll}
%     \displaystyle\sum_{j=1}^Jp_{ij}=w_i,&i=1,\ldots,J,\\[1mm]
%     \displaystyle\sum_{i=1}^J\sum_{j<2k}p_{ij}\le u_k,&k=1,\ldots,K,\\[1mm]
%     \displaystyle\sum_{i=1}^J\sum_{j\le2k}p_{ij}\ge\ell_k,&k=1,\ldots,K
%   \end{array}\right\}.
%   \label{eq-cell-mass-polytope}
% \end{equation}

% \begin{lemma}[Exact aggregation]
%   \label{lem-exact-aggregation}
%   For $\Xi=\mathbb R$,
%   \begin{equation}
%     \overline p_{\hat f}=\max_{p\in\mathcal A_\theta}\sum_{i,j=1}^J\overline h_{ij}p_{ij},
%     \qquad
%     \underline p_{\hat f}=\min_{p\in\mathcal A_\theta}\sum_{i,j=1}^J\underline h_{ij}p_{ij}.
%     \label{eq-finite-primal-values}
%   \end{equation}
% \end{lemma}
% \begin{proof}
%   Every feasible joint measure induces a matrix in $\mathcal A_\theta$. Its failure indicator on each cell pair lies between the coefficients in \eqref{eq-cell-costs}, giving one direction of each equality.

%   Conversely, take $p\in\mathcal A_\theta$. For $\bm x\in A_i$ with $w_i>0$, assign the residual to cell $C_j$ with probability $p_{ij}/w_i$. When $i\ne j$, any fixed representative of $C_j$ gives the prescribed indicator. The diagonal singleton uses its unique point. On a diagonal open interval, choosing $\xi=T(\bm x)$ attains $\overline h_{ii}=1$. To attain $\underline h_{ii}=0$, choose the midpoint of $T(\bm x)$ and the finite right endpoint, or $T(\bm x)+1$ on the last interval. These selections are measurable and stay in $C_j$. They define joint measures with input marginal $\mathbb P$, the cell masses $p_{ij}$, and the respective objectives in \eqref{eq-finite-primal-values}. Zero-mass input cells contribute nothing.
% \end{proof}

% On $C_j$, the dual step function has the constant value
% \begin{equation}
%   \eta_j:=
%   \sum_{k=1}^K\lambda_k\mathbf1_{\{j<2k\}}
%   -\sum_{k=1}^K\mu_k\mathbf1_{\{j\le2k\}}.
%   \label{eq-dual-cell-values}
% \end{equation}

% \begin{theorem}[Finite dual reformulations]
%   \label{thm-finite-dual-reformulations}
%   For $\Xi=\mathbb R$, both measure programs satisfy strong duality:
%   $\overline p_{\hat f}=\overline d_\theta$ and
%   $\underline p_{\hat f}=\underline d_\theta$. Their values are given by the finite linear programs
%   \begin{equation}
%     \begin{aligned}
%       \overline p_{\hat f}=\min_{\lambda,\mu\ge0,\ z\in\mathbb R^J}\quad&
%       \lambda^\top u-\mu^\top\ell+\sum_{i=1}^Jw_i z_i\\
%       \text{s.t.}\quad&
%       z_i+\sum_{k=1}^K\lambda_k\mathbf1_{\{j<2k\}}
%       -\sum_{k=1}^K\mu_k\mathbf1_{\{j\le2k\}}
%       \ge\overline h_{ij},
%       \quad i,j=1,\ldots,J,
%     \end{aligned}
%     \label{eq-finite-upper-dual}
%   \end{equation}
%   and
%   \begin{equation}
%     \begin{aligned}
%       \underline p_{\hat f}=\max_{\lambda,\mu\ge0,\ z\in\mathbb R^J}\quad&
%       -\lambda^\top u+\mu^\top\ell+\sum_{i=1}^Jw_i z_i\\
%       \text{s.t.}\quad&
%       z_i-\sum_{k=1}^K\lambda_k\mathbf1_{\{j<2k\}}
%       +\sum_{k=1}^K\mu_k\mathbf1_{\{j\le2k\}}
%       \le\underline h_{ij},
%       \quad i,j=1,\ldots,J.
%     \end{aligned}
%     \label{eq-finite-lower-dual}
%   \end{equation}
% \end{theorem}
% \begin{proof}
%   Dualizing the upper finite primal in \eqref{eq-finite-primal-values} gives \eqref{eq-finite-upper-dual}: $z_i$ is free for the row equality, and $\lambda_k,\mu_k\ge0$ correspond to the two prefix inequalities. Dualizing the lower finite primal gives \eqref{eq-finite-lower-dual}, with the indicated signs.

%   The polytope $\mathcal A_\theta$ is nonempty and compact. Nonemptiness follows from the product of the input law and the empirical residual law; compactness follows from the fixed row sums and nonnegativity. Both finite primal objectives are bounded between zero and one. Finite-dimensional LP strong duality therefore equates their values with their dual optima \citep{B.V04}.

%   It remains to identify these finite duals with the measure duals. For fixed multipliers and $\bm x\in A_i$, the upper pointwise constraints are equivalent to
%   $\varphi(\bm x)\ge\max_j(\overline h_{ij}-\eta_j)$.
%   Their least feasible potential is constant on each $A_i$ and can be written as $z_i$. This gives \eqref{eq-finite-upper-dual} exactly. Similarly, the lower constraints are equivalent to
%   $\varphi(\bm x)\le\min_j(\underline h_{ij}+\eta_j)$,
%   yielding \eqref{eq-finite-lower-dual}. The resulting potentials are bounded and measurable. Together with Lemma~\ref{lem-exact-aggregation}, these identities establish strong duality for the measure programs.
% \end{proof}

% Each dual has $4K+1$ scalar variables and $(2K+1)^2$ pointwise constraints. The cell probabilities $w_i$ are evaluated under $\mathbb P$; the aggregation preserves the original input law.

% \subsection{Closed-form bounds}
% \label{subsec-closed-form}

% The nested prefix constraints permit a further simplification of the finite LPs. Define
% \begin{equation}
%   H(t):=\mathbb P\{\hat g(\bm X)\le t\},
%   \qquad H_k:=H(-r_k)=\mathbb P\{T(\bm X)\ge r_k\}.
%   \label{eq-surrogate-response-cdf}
% \end{equation}
% For fixed residual cell masses $v_j\ge0$ with $\sum_jv_j=1$, let
% \[
%   V_m:=\sum_{j=1}^m v_j,\qquad
%   W_m:=\sum_{i=1}^m w_i,\qquad
%   V_0=W_0=0.
% \]
% Write $\Gamma(w,v)$ for the nonnegative matrices with row sums $w$ and column sums $v$. In particular, $H_k=1-W_{2k-1}$.

% \begin{lemma}[Ordered transport dual values]
%   \label{lem-ordered-transport}
%   For fixed $w$ and $v$,
%   \begin{align}
%     \max_{p\in\Gamma(w,v)}\sum_{i,j=1}^J\overline h_{ij}p_{ij}
%     &=\min_{m=0,\ldots,J}\{1-W_m+V_m\},
%     \label{eq-upper-transport-cuts}\\
%     \min_{p\in\Gamma(w,v)}\sum_{i,j=1}^J\underline h_{ij}p_{ij}
%     &=\max\left\{0,\max_{k=1,\ldots,K}(V_{2k}-W_{2k-1})\right\}.
%     \label{eq-lower-transport-cuts}
%   \end{align}
% \end{lemma}
% \begin{proof}
%   For the upper problem, form a bipartite flow network. Source-to-row arcs have capacities $w_i$, column-to-sink arcs have capacities $v_j$, and a row-to-column arc of capacity two is present exactly when $j\le i$. Every transport plan restricts to a flow on these permitted pairs. Conversely, any such flow can be completed to a transport plan using the remaining row and column masses. The maximal permitted mass therefore equals the maximum flow.

%   By the max-flow min-cut theorem \citep{F.F56}, this value equals the minimum cut capacity. A cut of capacity at most one exists, so a minimum cut crosses none of the capacity-two arcs. If the largest row index on the source side is $m$, every column $j\le m$ must also be on that side. Adding all rows $i\le m$ and removing columns $j>m$ from the source side cannot increase the cut capacity. Thus a minimum cut has capacity $1-W_m+V_m$. The empty row set gives $m=0$, proving \eqref{eq-upper-transport-cuts}.

%   For the lower problem, maximize the complementary safe mass. The permitted pairs are $j>i$ and the diagonal pairs with $i$ odd. The column neighborhoods are nested suffixes. For a nonempty source-side row set, its earliest permitted column is either one or $2k+1$ for some $k=1,\ldots,K$. In the latter case, replacing the source-side sets by all rows $i\ge2k$ and columns $j\ge2k+1$ cannot increase the cut capacity. The resulting capacity is
%   $W_{2k-1}+1-V_{2k}$. The empty row set and the suffix starting at column one both give capacity one. The same min-cut argument shows that the maximal safe mass is
%   $\min\{1,\min_k(W_{2k-1}+1-V_{2k})\}$.
%   Subtracting this value from one proves \eqref{eq-lower-transport-cuts}.
% \end{proof}

% \begin{theorem}[Closed-form failure probability bounds]
%   \label{thm-closed-form-bounds}
%   For $\Xi=\mathbb R$ and $\mathscr Q=\mathscr Q^{\mathrm{DW}}(\theta)$,
%   \begin{align}
%     \overline p_{\hat f}
%     &=\min\left\{1,\min_{k=1,\ldots,K}
%     \left[\beta_{c_{k-1}}(\theta)+H(-r_k)\right]\right\},
%     \label{eq-dw-closed-upper}\\
%     \underline p_{\hat f}
%     &=\max\left\{0,\max_{k=1,\ldots,K}
%     \left[\alpha_{c_k}(\theta)+H(-r_k)-1\right]\right\}.
%     \label{eq-dw-closed-lower}
%   \end{align}
% \end{theorem}
% \begin{proof}
%   The feasible residual prefix masses satisfy
%   \begin{equation}
%     0=V_0\le V_1\le\cdots\le V_J=1,\qquad
%     V_{2k-1}\le u_k,\quad V_{2k}\ge\ell_k.
%     \label{eq-feasible-prefix-masses}
%   \end{equation}
%   Conversely, any sequence satisfying \eqref{eq-feasible-prefix-masses} defines column masses $v_j=V_j-V_{j-1}$ and admits a transport plan with row sums $w$, for example $p_{ij}=w_i v_j$.

%   Set $\ell_0:=0$ and $u_{K+1}:=1$. The rank endpoints imply that $u_k$ and $\ell_k$ are nondecreasing, with $\ell_{k-1}\le u_k$ and $\ell_k\le u_{k+1}$. Define two feasible prefix sequences by
%   \begin{equation}
%     \begin{aligned}
%       V^\uparrow_{2k-1}&:=u_k,&
%       V^\uparrow_{2k}&:=u_{k+1},\\
%       V^\downarrow_{2k-1}&:=\ell_{k-1},&
%       V^\downarrow_{2k}&:=\ell_k,
%       \qquad k=1,\ldots,K,
%     \end{aligned}
%     \label{eq-extreme-prefix-vectors}
%   \end{equation}
%   with $V^\uparrow_0=V^\downarrow_0=0$ and
%   $V^\uparrow_J=V^\downarrow_J=1$.
%   Every feasible sequence satisfies
%   $V^\downarrow_m\le V_m\le V^\uparrow_m$:
%   an even prefix is bounded above by the following odd prefix, and an odd prefix is bounded below by the preceding even prefix.

%   The right-hand side of \eqref{eq-upper-transport-cuts} is nondecreasing in every $V_m$, so the upper problem is optimized at $V^\uparrow$. Its odd-index cuts have values
%   \[
%     1-W_{2k-1}+V^\uparrow_{2k-1}=H_k+u_k.
%   \]
%   For $k<K$, the cut at $m=2k$ is dominated by that at $m=2k+1$, since both have prefix value $u_{k+1}$ and $W_m$ is nondecreasing. The cut at $m=2K$ is at least one, while $m=0$ and $m=J$ give one. This proves \eqref{eq-dw-closed-upper}.

%   The right-hand side of \eqref{eq-lower-transport-cuts} is nondecreasing in the even prefixes, so its minimum occurs at $V^\downarrow$. Substitution gives
%   $\max\{0,\max_k(\ell_k+H_k-1)\}$, proving \eqref{eq-dw-closed-lower}.
% \end{proof}

% The cut classification also identifies optimal solutions of the measure duals. If $u_k+H_k=\overline p_{\hat f}$, one optimal certificate is
% \[
%   \lambda=e_k,\qquad \mu=0,\qquad
%   \varphi(\bm x)=\mathbf1_{\{\hat g(\bm x)\le-r_k\}},
% \]
% where $e_k$ is the $k$th unit vector. When $\overline p_{\hat f}=1$, an optimal certificate is $\lambda=\mu=0$ and $\varphi\equiv1$. If $\ell_k+H_k-1=\underline p_{\hat f}$, an optimal lower certificate is
% \[
%   \lambda=0,\qquad \mu=e_k,\qquad
%   \varphi(\bm x)=\mathbf1_{\{\hat g(\bm x)\le-r_k\}}-1;
% \]
% when $\underline p_{\hat f}=0$, an optimal certificate uses $\lambda=\mu=0$ and $\varphi\equiv0$. Thus an optimal dual solution can be chosen with at most one positive band multiplier. Its optimality follows from the finite LP and min-cut derivation above.

% \begin{remark}[Known interval support]
%   \label{rem-interval-support}
%   If $\Xi=[a,b]$ is known in advance, the measure programs and their duals retain the domains in \eqref{eq-measure-upper}--\eqref{eq-lower-measure-dual}. Adding $a$ and $b$ to the ordered partition and removing residual cells outside $[a,b]$ gives the corresponding finite LPs. Input thresholds below $a$ cannot fail, and thresholds at least $b$ must fail. The endpoint cuts consequently replace the outer constants one and zero in \eqref{eq-dw-closed-upper}--\eqref{eq-dw-closed-lower} by $H(-a)$ and $H(-b)$, respectively.
% \end{remark}

% After rank endpoint inversion, \eqref{eq-dw-closed-upper}--\eqref{eq-dw-closed-lower} require only $K$ evaluations of the surrogate response CDF and finite minimum and maximum operations. The upper candidate $u_k+H_k$ combines the permitted residual left-tail mass with the input probability whose surrogate margin is insufficient at that threshold. The lower candidate $\ell_k+H_k-1$ measures unavoidable overlap of the corresponding events.

% For $\theta\ge\theta_{N^{\mathrm{te}}}(\delta)$, these exact reformulations inherit the coverage guarantee in Theorem~\ref{theorem-01}. If the CDF values are evaluated by simulation, simultaneous upper confidence bounds for $H_k$ yield a certified upper failure bound, and simultaneous lower confidence bounds yield a certified lower bound. The numerical integration error is then included in the overall confidence budget.

\bibliographystyle{elsarticle-harv}
\bibliography{reference}

\end{document}
